\chapter{Evaluation}

\section{Evaluation metrics}

	\subsection{Retriever Metrics}
	
	\textbf{Recall@K:}  
	Recall@K quantifies the proportion of relevant documents successfully retrieved among the top $K$ results for a given query.
	\[
	\text{Recall@K} = \frac{| \text{Relevant documents} \cap \text{Top-}K \text{ retrieved documents} |}{|\text{Relevant documents}|}
	\]
	Where:
	\begin{itemize}
		\item $|\text{Relevant documents}|$ denotes the total number of documents relevant to the query.
		\item $|\text{Relevant documents} \cap \text{Top-}K \text{ retrieved documents}|$ is the count of relevant documents found in the top $K$ results.
	\end{itemize}
	\textbf{Example:}
	\begin{itemize}
		\item Scenario: A query returns the top \( K = 5 \) documents. Among them, \( 3 \) are relevant, while the total number of relevant documents is \( 4 \).
		\item Calculation:
		\begin{equation*}
			\text{Recall@5} = \frac{3}{4} = 0.75 \quad (75\%)
		\end{equation*}
	\end{itemize}
	
	\textbf{Precision@K:}  
	Precision@K measures the accuracy of the top $K$ retrieved documents by evaluating how many of them are truly relevant.
	\[
	\text{Precision@K} = \frac{| \text{Relevant documents} \cap \text{Top-}K \text{ retrieved documents} |}{|\text{Top-}K \text{ retrieved documents}|}
	\]
	Where:
	\begin{itemize}
		\item $|\text{Top-}K \text{ retrieved documents}|$ is the number of documents retrieved within the top $K$ results.
	\end{itemize}
	\textbf{Example:}
	\begin{itemize}
		\item Scenario: A query retrieves the top \( K = 5 \) documents, with \( 3 \) being relevant. The total number of relevant documents is \( 4 \).
		\item Calculation:
		\begin{equation*}
			\text{Precision@5} = \frac{3}{5} = 0.6 \quad (60\%)
		\end{equation*}
	\end{itemize}
	
	\textbf{Latency:}  
	Latency represents the response time taken to complete a retrieval operation.
	\[
	\text{Latency (ms)} = \text{Time}_{\text{returned}} - \text{Time}_{\text{submitted}}
	\]
	\textbf{Example:}
	\begin{itemize}
		\item Scenario: The query is submitted at 11:00:00.000 and the result is received at 11:00:00.340.
		\item Calculation:
		\begin{equation*}
			\text{Latency} = 340\,\text{ms} - 0\,\text{ms} = 340\,\text{ms}
		\end{equation*}
	\end{itemize}
	
	\subsection{Generator Metrics}
	
	\textbf{ROUGE Score:}  
	ROUGE evaluates the lexical overlap between generated content and a set of reference texts. ROUGE-N specifically computes $n$-gram overlaps:
	\[
	\text{ROUGE-N} = \frac{\sum_{\text{ref} \in \text{References}} \sum_{\text{ngram} \in \text{Generated}} \text{Count}(\text{ngram})}{\sum_{\text{ref} \in \text{References}} \sum_{\text{ngram} \in \text{ref}} \text{Count}(\text{ngram})}
	\]
	\textbf{Example:}
	\begin{itemize}
		\item Scenario: A reference document has 10 bigrams, and the generated output contains 8 matching bigrams.
		\item Calculation:
		\begin{equation*}
			\text{ROUGE-2} = \frac{8}{10} = 0.8
		\end{equation*}
	\end{itemize}
	
	\textbf{Relevance:}  
	Relevance measures how well the generated outputs align with the intent of user queries, typically rated manually and averaged across queries.
	\[
	\text{Relevance} = \frac{1}{N} \sum_{i=1}^{N} \text{RelevanceScore}_i
	\]
	\textbf{Example:}
	\begin{itemize}
		\item Scenario: Five queries yield relevance scores of \( 0.9, 0.8, 0.7, 1.0, \text{and } 0.6 \).
		\item Calculation:
		\begin{equation*}
			\text{Relevance} = \frac{0.9 + 0.8 + 0.7 + 1.0 + 0.6}{5} = 0.8
		\end{equation*}
	\end{itemize}
	
	\textbf{Latency:}  
	Latency in the generation phase denotes the time elapsed between the completion of retrieval and the generation of the final output.
	\[
	\text{Latency (ms)} = \text{Time}_{\text{generated}} - \text{Time}_{\text{retrieved}}
	\]
	\textbf{Example:}
	\begin{itemize}
		\item Scenario: The output is generated at 12:00:00.210 after retrieval completes at 12:00:00.000.
		\item Calculation:
		\begin{equation*}
			\text{Latency} = 210\,\text{ms} - 0\,\text{ms} = 210\,\text{ms}
		\end{equation*}
	\end{itemize}


\section{Evaluation methodology}

To evaluate the performance of both the retriever and the generator, we employ a curated dataset containing queries with predefined relevant documents. The evaluation covers two types:

\begin{itemize} 
	\item \textbf{Simple queries:} Direct questions that require retrieving clearly defined information from the documentation. 
	\item \textbf{Complex queries:} Queries involving multiple components, conditions, or requiring deeper semantic understanding.
\end{itemize}

\section{Evaluation summary}

\subsection{Retriever Performance}

Table \ref{tab:ret_simple} and Table \ref{tab:ret_complex} present the retriever performance and generator latency for simple and complex queries, respectively.

Results show that the retriever achieved high \textbf{Recall@K} in simple queries, with a peak of \textbf{0.81} at \textbf{K=5}, demonstrating strong coverage of relevant documents. However, performance declined in complex queries due to their nuanced requirements, with \textbf{Recall@K} dropping to \textbf{0.24} at \textbf{K=5}. \textbf{Precision@K} also showed a similar pattern, peaking at \textbf{0.57} in simple cases but falling below \textbf{0.3} for complex ones.

Latency remained reasonable for simple queries (approximately \textbf{1.2 seconds}), while complex queries understandably took longer due to the increased search depth and token length, averaging around \textbf{3.4 seconds}.


\subsection{Generator Performance}

The generator ROUGE scores across various models are illustrated in Figure~\ref{fig:rouge_1}, Figure~\ref{fig:rouge_2} and Figure~\ref{fig:rouge_3}, corresponding to ROUGE-1, ROUGE-2 and ROUGE-3 metrics, respectively. Additionally, model-wise relevance evaluations are presented in Figure~\ref{fig:llm_eval_small} and Figure~\ref{fig:llm_eval_bigger}, capturing the perceived quality of generated outputs across different model sizes. Notably:

\begin{itemize}
	\item \textbf{Codestral-22B}, \textbf{LLaMA2-13B}, and \textbf{Qwen2.5-14B} consistently achieved top scores, producing highly relevant and fluent outputs across diverse test cases.
	\item \textbf{Mistral-7B}, \textbf{CodeLLaMA-7B}, and \textbf{Qwen2.5-3B} also demonstrated balanced performance, with strong semantic accuracy and reliable response generation.
	\item Smaller models like \textbf{Qwen-1.8B} and \textbf{CodeGemma-2B} showed limitations in deeper contextual understanding, reflected in lower performance in complex reasoning tasks.
\end{itemize}

Overall, generation latency remained within practical limits—averaging approximately \textbf{6 seconds} for simple prompts and \textbf{12 seconds} for complex queries—though these values varied depending on the size and architecture of the underlying language model.

\section{Testing}

\subsection{Unit testing} Validate individual functions and components of the extension.

\subsection{Integration testing} Ensure proper interaction between extension components and the VSCode environment.